% ----------------------------------------------------------
% ELEMENTOS PRÉ-TEXTUAIS
% ----------------------------------------------------------
% \pretextual

% ---
% Capa
% ---
\imprimircapa
% ---

% ---
% Folha de rosto
% (o * indica que haverá a ficha bibliográfica)
% ---
%\imprimirfolhaderosto*
\imprimirfolhaderosto
% ---



% ---
% Agradecimentos
% ---
%\begin{agradecimentos}
%Não remover a palavra agradecimentos.
%Os agradecimentos devem ser dirigidos a quem realmente contribuiu, de maneira relevante, à elaboração do trabalho. 

%\end{agradecimentos}
% ---

% ---
% Epígrafe
% ---
%\begin{epigrafe}
%    \vspace*{\fill}
%	\begin{flushright}
%		\textit{Epígrafe (Não se escreve a palavra epígrafe). Elemento opcional. 
%A epígrafe deve ser colocada após o agradecimento; trata-se de uma citação, seguida de indicação de autoria, relacionada à matéria tratada no corpo do trabalho. Deve ser “[...] elaborada conforme a NBR 10520 [...]. Podem também constar epígrafes nas folhas ou páginas de abertura das seções primárias” (ABNT, 2002, p. 7). 
%A fonte da epígrafe deve sempre ser mencionada nas referências. 
%Citação direta até 3 linhas deve estar entre aspas e em parágrafo normal (vá até a janela de Estilo - selecione - Parágrafo), se tiver mais de 3 linhas, deve ser recuada 4 cm da margem esquerda, com fonte menor que 12 e espaçamento entre linhas simples 
%}
%	\end{flushright}
%\end{epigrafe}
% ---

% ---
% RESUMOS
% ---

% resumo em português
\setlength{\absparsep}{18pt} % ajusta o espaçamento dos parágrafos do resumo
\begin{resumo}
Os transformadores de potência são equipamentos essenciais nos sistemas de energia elétrica, apresentando alto custo e grande impacto operacional. Falhas nesses equipamentos podem provocar interrupções no fornecimento de energia e prejuízos significativos. Nesse contexto, a Análise de Gases Dissolvidos (DGA) destaca-se como uma das principais técnicas para o monitoramento da condição interna dos transformadores. Entretanto, métodos tradicionais de interpretação dos resultados da DGA, como o Triângulo de Duval e as razões da norma IEC 60599, apresentam limitações devido às suas fronteiras fixas, o que pode levar a diagnósticos imprecisos ou indeterminados. Com o avanço das técnicas de Inteligência Artificial, novas abordagens têm sido propostas para superar essas limitações. No entanto, grande parte dessas soluções ainda se concentra na identificação de falhas já existentes, caracterizando uma abordagem reativa. Diante disso, este trabalho desenvolve e avalia um sistema híbrido baseado em redes neurais recorrentes LSTM para previsão de gases e no classificador XGBoost otimizado por Algoritmo Genético Melhorado (IGA) para o diagnóstico de falhas. A investigação foi estruturada em dois eixos experimentais para contemplar as características dos dados industriais abertos disponíveis. No eixo A, foram analisadas 3.000 séries temporais com quatro gases para avaliar se a inclusão das projeções da LSTM melhora o desempenho do classificador. Os testes pareados por bootstrap demonstraram que, devido à forte correlação entre as leituras do momento e as projeções futuras, as estimativas temporais não trouxeram ganho estatístico no diagnóstico pontual. Por outro lado, o algoritmo genético selecionou um subconjunto de apenas 8 atributos a partir dos 17 candidatos, mantendo o mesmo desempenho ($F_1 \approx 0,89$). No eixo B, com 1.972 registros estáticos de cinco gases e sete classes de falhas, os modelos baseados em árvores alcançaram cerca de 76\% de $F_1$-macro, sugerindo que a sobreposição física entre diferentes tipos de degradação impõe um limite prático ao diagnóstico baseado unicamente em gases. Por fim, a técnica TreeSHAP permitiu analisar as regras aprendidas pelo modelo, evidenciando boa coerência entre as fronteiras do XGBoost e as faixas térmicas e elétricas das normas clássicas, além de explicar os casos fronteiriços de erro de classificação.

 \textbf{Palavras-chave}: Transformadores de Potência, Análise de Gases Dissolvidos, Redes LSTM, XGBoost, Algoritmo Genético Melhorado, TreeSHAP.
\end{resumo}

% resumo em inglês
\begin{resumo}[Abstract]
 \begin{otherlanguage*}{english}
Power transformers are essential components in electrical power systems, involving high costs and significant operational impact. Failures in these devices can lead to power supply interruptions and substantial economic losses. In this context, Dissolved Gas Analysis (DGA) stands out as one of the main techniques for monitoring the internal condition of transformers. However, traditional methods for interpreting DGA results, such as the Duval Triangle and Rogers Ratio Method, present limitations due to their fixed boundaries, which may result in imprecise or inconclusive diagnoses. With the advancement of Artificial Intelligence techniques, new approaches have been proposed to overcome these limitations. Nevertheless, most of these solutions still focus on identifying faults that have already occurred, representing a reactive approach. Therefore, this work develops and evaluates a system based on Long Short-Term Memory (LSTM) recurrent neural networks for gas forecasting and an XGBoost classifier optimized by an Improved Genetic Algorithm (IGA) for fault diagnosis. The investigation was structured into two experimental tracks to accommodate the characteristics of available open industrial datasets. In Track A, 3,000 time series covering four gases were analyzed to test whether adding LSTM forecasts improves classification performance. Paired bootstrap tests showed that, due to the high correlation between current concentrations and future projections, temporal estimates provided no statistically significant gain for instantaneous diagnosis; on the other hand, the genetic algorithm reduced the required feature set from 17 to 8 variables while maintaining the same diagnostic performance ($F_1 \approx 0.89$). In Track B, using 1,972 static records with five gases and seven fault classes, tree-based models reached approximately 76\% in $F_1$-macro, indicating that physical overlap between fault mechanisms imposes a practical limit on pure gas-based diagnosis. Finally, TreeSHAP was applied to explain the learned decision patterns, showing strong alignment between XGBoost boundaries and traditional standard thresholds, while clarifying borderline classification errors.

   \vspace{\onelineskip}
 
   \noindent 
   \textbf{Keywords}: Power Transformers, Dissolved Gas Analysis, LSTM Networks, XGBoost, Improved Genetic Algorithm, TreeSHAP.
 \end{otherlanguage*}
\end{resumo}


% ---
% inserir lista de ilustrações
% ---
\pdfbookmark[0]{\listfigurename}{lof}
\listoffigures*
\cleardoublepage
% ---

% ---
% inserir lista de quadros
% ---
%\pdfbookmark[0]{\listofquadrosname}{loq}
¨%\listofquadros*
%\cleardoublepage
% ---

% ---
% inserir lista de tabelas
% ---
\pdfbookmark[0]{\listtablename}{lot}
\listoftables*
\cleardoublepage
% ---

% ---
% inserir lista de abreviaturas e siglas
% ---
\begin{siglas}
  \item[ABNT] Associação Brasileira de Normas Técnicas
  \item[AG] Algoritmo Genético
  \item[ANFIS] \textit{Adaptive Neuro-Fuzzy Inference System} (Sistema de Inferência Neuro-Fuzzy Adaptativo)
  \item[BPTT] \textit{Backpropagation Through Time} (Retropropagação Através do Tempo)
  \item[CNN] \textit{Convolutional Neural Network} (Rede Neural Convolucional)
  \item[CRISP-DM] \textit{Cross-Industry Standard Process for Data Mining} (Processo Padrão Intersetorial para Mineração de Dados)
  \item[DGA] \textit{Dissolved Gas Analysis} (Análise de Gases Dissolvidos)
  \item[GA] \textit{Genetic Algorithm} (Algoritmo Genético)
  \item[IEC] \textit{International Electrotechnical Commission} (Comissão Eletrotécnica Internacional)
  \item[IGA] \textit{Improved Genetic Algorithm} (Algoritmo Genético Melhorado)
  \item[KNN] \textit{k-Nearest Neighbors} (k-Vizinhos Mais Próximos)
  \item[LightGBM] \textit{Light Gradient Boosting Machine} (Máquina de Aumento de Gradiente Leve)
  \item[LSTM] \textit{Long Short-Term Memory} (Memória de Curto e Longo Prazo)
  \item[MAE] \textit{Mean Absolute Error} (Erro Médio Absoluto)
  \item[MLP] \textit{Multilayer Perceptron} (Perceptron Multicamadas)
  \item[PSO] \textit{Particle Swarm Optimization} (Otimização por Enxame de Partículas)
  \item[RMSE] \textit{Root Mean Squared Error} (Raiz do Erro Quadrático Médio)
  \item[RNN] \textit{Recurrent Neural Network} (Rede Neural Recorrente)
  \item[SHAP] \textit{SHapley Additive exPlanations} (Explicações Aditivas de Shapley)
  \item[SMOTE] \textit{Synthetic Minority Over-sampling Technique} (Técnica de Sobreamostragem de Minorias Sintéticas)
  \item[SVM] \textit{Support Vector Machine} (Máquina de Vetores de Suporte)
  \item[XGBoost] \textit{Extreme Gradient Boosting} (Aumento Extremo de Gradiente)
\end{siglas}
% ---

% ---
% inserir lista de símbolos
% ---
\begin{simbolos}
  \item[$w$] Tamanho da janela temporal (número de passos históricos)
  \item[$h$] Horizonte de previsão (passos à frente)
  \item[$x_t$] Vetor de concentrações dos gases no instante $t$
  \item[$\hat{y}$] Saída predita pelo modelo
  \item[$h_t$] Estado oculto da rede recorrente no instante $t$
  \item[$C_t$] Estado da célula (memória) da LSTM no instante $t$
  \item[$f_k(x)$] Árvore $k$ no ensemble de \textit{gradient boosting}
  \item[$L(y,\hat{y})$] Função de perda
  \item[$N$] Tamanho da população no AG/IGA
  \item[$G$] Número máximo de gerações no AG/IGA
  \item[$P_c$] Taxa de cruzamento (\textit{crossover rate})
  \item[$P_m$] Taxa de mutação
\end{simbolos}
% ---

% ---
% inserir o sumario
% ---
\pdfbookmark[0]{\contentsname}{toc}
\tableofcontents*
\cleardoublepage
% ---